\documentclass[12pt]{article}
\usepackage[T1]{fontenc}
\usepackage[utf8]{inputenc}
\usepackage{lmodern}
\usepackage{textcomp}
\usepackage{enumitem}
\usepackage{geometry}
\geometry{margin=1in}
\begin{document}
\begin{center}
Answer Key - Psychology - Graduate Level Assessment 16 \
\small (Answers are ordered exactly as in the question paper.)
\end{center}

\section*{Multiple Choice}
\begin{enumerate}[label=\arabic*.]
\item \textbf{Answer:} A
\\ \textit{Options:} A) "Genetic blueprint" hypothesis; B) "Innate potential" hypothesis; C) Environmental impact hypothesis; D) "Experience-driven" development theory; E) "Blank slate" hypothesis
\\ \textit{Note:} The "genetic blueprint" hypothesis proposed by Stern in 1956 emphasizes the significant role of genetic inheritance in shaping individual characteristics, thereby positioning it as a key argument in the nature-nurture controversy.
\item \textbf{Answer:} C
\\ \textit{Options:} A) too little curvature of the cornea and lens; B) too much curvature of the retina and lens; C) too much curvature of the cornea and lens; D) too much curvature of the iris and cornea; E) incorrect alignment of the cornea and retina
\\ \textit{Note:} Nearsightedness, or myopia, occurs when the eye's cornea and lens are excessively curved, causing light rays to focus in front of the retina rather than directly on it, leading to blurred distance vision.
\item \textbf{Answer:} A
\\ \textit{Options:} A) By observing individual's behaviors and actions; B) Through the analysis of social media usage and engagement patterns; C) By analyzing the country's economic factors; D) By examining the literacy and numeracy rates; E) Through standardized achievement tests in educational institutions
\\ \textit{Note:} Achievement Motive (nAch) is assessed through the observation of behaviors and actions because these manifestations reveal individuals' drive for success and their approach to challenges, which can also be extrapolated to understand cultural or national trends in achievement motivation.
\item \textbf{Answer:} A
\\ \textit{Options:} A) Recovery rate in individuals using self-help resources; B) Incidence in the population of the disorder being eared; C) Recovery rate in patients using alternative therapies; D) Improvement rate in patients receiving no therapy; E) Cure rate for older, outdated treatment methods
\\ \textit{Note:} The recovery rate in individuals using self-help resources serves as a proper comparison rate because it provides a baseline for evaluating the effectiveness of the new psychotherapy against an established method of treatment that emphasizes self-directed healing.
\item \textbf{Answer:} B
\\ \textit{Options:} A) Recognize faces; B) move his left hand; C) Hear sounds correctly; D) understand information he reads; E) speak intelligibly
\\ \textit{Note:} The correct answer is based on the principle of contralateral control in the brain, where damage to the right frontal lobe affects motor functions on the opposite side of the body, resulting in Mr. Gordon's inability to move his left hand.
\end{enumerate}

\section*{True / False}
\begin{enumerate}[label=\arabic*.]
\setcounter{enumi}{5}
\item \textbf{Answer:} False
\\ \textit{Options:} A) True; B) False
\\ \textit{Note:} Research in cognitive psychology indicates that while there are some differences in spatial reasoning abilities between genders, studies show that men often outperform women in spatial reasoning tasks, making the statement false.
\item \textbf{Answer:} True
\\ \textit{Options:} A) True; B) False
\\ \textit{Note:} In classical conditioning, the process involves the learner forming an association between the unconditioned stimulus, which naturally elicits a response, and the conditioned stimulus, which initially does not provoke that response but eventually does after repeated pairings.
\item \textbf{Answer:} False
\\ \textit{Options:} A) True; B) False
\\ \textit{Note:} Encouraging competitive sports and games can often heighten aggressiveness in children due to the inherent pressure to win and the potential for conflict in competitive environments, rather than reducing it.
\item \textbf{Answer:} True
\\ \textit{Options:} A) True; B) False
\\ \textit{Note:} The Draw-a-Person Test (DAP) evaluates fine motor coordination by analyzing the individual's ability to represent human figures through drawing, which requires the use of precise motor skills.
\item \textbf{Answer:} True
\\ \textit{Options:} A) True; B) False
\\ \textit{Note:} The statement is true because prolonged delays between a conditioned stimulus and an unconditioned stimulus can disrupt the temporal contiguity necessary for effective associative learning, making it less likely for the animal to form a strong connection between the two stimuli.
\end{enumerate}

\section*{Long Form Response}
\begin{enumerate}[label=\arabic*.]
\setcounter{enumi}{10}
\item \textbf{Answer:} Garcia and Koelling's research on learned aversions significantly highlights the role of pleasure as an unconditioned stimulus (UCS) when paired with a visual stimulus as a conditioned stimulus (CS). Their findings indicated that organisms can develop aversions to a taste (the UCS) when paired with a specific visual cue (the CS), demonstrating that the association is robust even after a delay. However, when pleasure was employed as a UCS alongside a visual stimulus, the results suggested that positive experiences could also create associative learning, albeit less frequently studied. This research underscores the complexity of conditioning processes, revealing that while negative associations are often more potent, positive stimuli can still influence behavior when contextually appropriate, thus broadening our understanding of the mechanisms behind learned aversions.
\\ \textit{Note:} correct because it accurately reflects Garcia and Koelling's demonstration of how learned aversions can be established through the pairing of an unconditioned stimulus (like taste) with a conditioned stimulus (such as a visual cue), while also indicating that the effectiveness of pleasure as a UCS in such associations is significant, particularly in the context of delayed conditioning.
\item \textbf{Answer:} The anxiety theory of neurosis posits that neurotic symptoms often stem from unresolved conflicts and anxiety rooted in childhood trauma. Such traumas, which may include neglect, abuse, or loss, can disrupt normal emotional development and lead to maladaptive coping mechanisms. As individuals navigate their adult lives, these unresolved childhood experiences may manifest as anxiety disorders, phobias, or obsessive-compulsive behaviors. Understanding the link between childhood trauma and neurotic symptoms is crucial for psychological treatment, as it underscores the importance of addressing past experiences in therapeutic settings. Effective treatment approaches, such as psychodynamic therapy or trauma-informed care, focus on processing these early experiences to alleviate symptoms and promote healthier coping strategies.
\\ \textit{Note:} correct because it accurately describes how unresolved childhood traumas can disrupt emotional development, leading to neurotic symptoms and highlighting the importance of addressing these issues in psychological treatment.
\end{enumerate}

\vfill
\noindent\textit{End of Answer Key}
\end{document}